%% 35장 수치 적분 I
\chapter{수치 적분 I}
\label{ch:quadrature-1}

\abstract{정적분을 함수의 값 몇 개의 가중합으로 어림하는 규칙을 구적 규칙이라 한다. 이 장에서는 보간 다항식을 적분하는 Newton--Cotes 규칙에서 시작해 고른 간격의 노드에서 차수를 올리면 무게가 음수가 되어 규칙이 불안정해지는 것을 정확한 유리수 계산으로 본다. 복합 사다리꼴 규칙의 오차는 Euler--Maclaurin 공식으로 이해되고, 주기 함수에서는 그 규칙이 기하급수적으로 수렴한다. 이어 직교다항식의 세 항 점화식과 근의 성질을 증명하고 노드를 직교다항식의 근으로 잡은 $n$점 Gauss 구적이 차수 $2n-1$까지 정확하고 무게가 양수임을 증명한다. Golub--Welsch 방법은 노드와 무게를 삼중대각 행렬의 고유값 문제로 구하지만 작은 무게에서 정확도를 잃는다. Chebyshev 점의 Clenshaw--Curtis 규칙은 Gauss 구적에 거의 뒤지지 않는다. 끝으로 오차를 어림하는 Gauss--Kronrod 쌍과 그 어림이 틀리는 경우를 본다.}

\section{보간형 구적}
\label{sec:interpolatory-quadrature}

\subsection{정확도의 차수}

$\int_a^bf(x)\,dx$를 $\sum_{i=0}^nw_if(x_i)$로 어림하는 것을 \term{구적 규칙}{quadrature rule}이라 하고 $x_i$를 노드, $w_i$를 무게라 한다. 차수 $d$ 이하의 모든 다항식을 정확히 적분하는 가장 큰 $d$를 규칙의 \term{정확도의 차수}{degree of exactness}라 한다. 노드가 주어지면 가장 자연스러운 무게는 29장의 보간 다항식을 적분한 것이다.
\[
w_i=\int_a^bL_i(x)\,dx
\]
이렇게 만든 \term{보간형 구적}{interpolatory quadrature}은 노드가 $n+1$개이면 차수 $n$ 이하의 다항식에서 정확하고, 그 반대도 성립한다. 차수 $n$까지 정확한 규칙은 $L_i$를 정확히 적분하므로 무게가 바로 $\int L_i$이기 때문이다.

구적 규칙의 안정성은 무게의 크기로 정해진다. 함숫값에 크기 $\delta$의 오차가 있으면 결과의 오차는 많아야 $\delta\sum_i|w_i|$이고, 상수 $1$을 정확히 적분하면 $\sum_iw_i=b-a$이다. 따라서 $\sum_i|w_i|/(b-a)$가 구적의 조건수 구실을 한다. 무게가 모두 양수이면 이 값은 정확히 $1$이고, 음수 무게가 있으면 $1$보다 커진다.

\subsection{Newton--Cotes 규칙}

노드를 고른 간격으로 두면 \term{Newton--Cotes 규칙}{Newton--Cotes rule}이다. 양 끝을 포함한 닫힌 규칙의 $N=1$은 사다리꼴 규칙, $N=2$는 Simpson 규칙이다. 무게는 정확한 유리수로 계산할 수 있는데, $[0,N]$의 정수 노드에서 결과를 보면 음수 무게가 처음 나타나는 것은 $N=8$이다. $N=8$에서 음수 무게가 $3$개 생기고 $\sum|w_i|/N=1.451$이다. $N=10$에서는 음수가 $4$개, $\sum|w_i|/N=3.065$이고, $N=12$에서는 $7.53$, $N=20$에서는 $544$이며 가장 작은 무게가 $-90N$이다. 29장의 Lebesgue 상수가 고른 간격에서 지수적으로 자라던 것과 같은 현상이다. Pólya는 연속인 함수는 물론 해석 함수 가운데에도 Newton--Cotes 규칙이 $N\to\infty$에서 수렴하지 않는 것이 있음을 보였다\cite{polya1933}. 그래서 Newton--Cotes 규칙은 높은 차수로 쓰지 않고, 구간을 나누어 낮은 차수의 규칙을 되풀이하는 복합 규칙으로 쓴다.

\section{복합 규칙}
\label{sec:composite}

\subsection{사다리꼴 규칙과 Euler--Maclaurin 공식}

구간을 길이 $h=(b-a)/n$인 $n$조각으로 나누고 조각마다 사다리꼴 규칙을 쓰면 \term{복합 사다리꼴 규칙}{composite trapezoidal rule}
\[
T_h=h\Bigl(\tfrac12f(x_0)+f(x_1)+\dots+f(x_{n-1})+\tfrac12f(x_n)\Bigr)
\]
을 얻는다. 그 오차는 \term{Euler--Maclaurin 공식}{Euler--Maclaurin formula}이 정확히 말해 준다. $f$가 충분히 매끄러우면
\begin{equation}
\label{eq:euler-maclaurin}
T_h-\int_a^bf=\sum_{k=1}^{m}\frac{B_{2k}h^{2k}}{(2k)!}\bigl(f^{(2k-1)}(b)-f^{(2k-1)}(a)\bigr)+O(h^{2m+2})
\end{equation}
이고, $B_{2k}$는 Bernoulli 수이다\cite{atkinson1989}. 오차가 $h^2$의 거듭제곱 급수로 전개되므로 33장의 Richardson 외삽을 그대로 쓸 수 있고, 이것이 \term{Romberg 적분}{Romberg integration}이다. $[0,1]$에서 $e^x$를 적분하면 $n=4,8,16,32$에서 오차가 $8.9\times10^{-3}$, $2.2\times10^{-3}$, $5.6\times10^{-4}$, $1.4\times10^{-4}$로 정확히 $O(h^2)$이다.

\subsection{주기 함수}

\eqref{eq:euler-maclaurin}의 모든 항은 양 끝의 도함수의 차에 비례한다. 그러므로 $f$가 주기 $b-a$인 매끄러운 주기 함수이면 모든 항이 사라지고, 오차는 $h$의 어떤 거듭제곱보다 빨리 준다. $f$가 실수축 둘레의 띠에서 해석적이면 오차는 기하급수적으로 준다\cite{trefethen2014trapezoidal}. $\int_0^{2\pi}e^{\cos x}\,dx=2\pi I_0(1)$을 사다리꼴 규칙으로 계산하면 $n=4$에서 상대오차가 $4.3\times10^{-3}$, $n=8$에서 $1.6\times10^{-7}$이고 $n=16$에서는 binary64로 정확하다. 가장 단순한 규칙이 주기 함수에서는 최선에 가깝다는 이 사실은 37장의 이산 푸리에 변환과 36장의 이중 지수 구적의 바탕이다.

\section{직교다항식}
\label{sec:orthogonal-polynomials}

\subsection{세 항 점화식}

$[a,b]$ 위의 양의 무게 함수 $\omega$에 대해 내적 $\langle f,g\rangle=\int_a^bf(x)g(x)\omega(x)\,dx$를 두고, 단항식 $1,x,x^2,\dots$를 이 내적으로 Gram--Schmidt 직교화하면 \term{직교다항식}{orthogonal polynomials} $p_0,p_1,\dots$을 얻는다. $p_k$는 차수가 정확히 $k$이고 차수가 $k$보다 낮은 모든 다항식과 직교한다. $\omega=1$, $[-1,1]$이면 Legendre 다항식, $\omega=(1-x^2)^{-1/2}$이면 Chebyshev 다항식이다.

\begin{theorem}
\label{thm:three-term}
최고차 계수가 $1$인 직교다항식은 $p_{-1}=0$, $p_0=1$과
\[
p_{k+1}(x)=(x-\alpha_k)p_k(x)-\beta_kp_{k-1}(x),\qquad\alpha_k=\frac{\langle xp_k,p_k\rangle}{\langle p_k,p_k\rangle},\quad\beta_k=\frac{\langle p_k,p_k\rangle}{\langle p_{k-1},p_{k-1}\rangle}
\]
를 만족한다.
\end{theorem}

\begin{proof}
$xp_k-p_{k+1}$은 차수가 $k$ 이하이므로 $p_0,\dots,p_k$의 일차결합 $\sum_{j\le k}c_jp_j$로 쓸 수 있고, $c_j=\langle xp_k,p_j\rangle/\langle p_j,p_j\rangle$이다. 그런데 $\langle xp_k,p_j\rangle=\langle p_k,xp_j\rangle$이고 $xp_j$의 차수가 $j+1$이므로, $j+1<k$이면 이 값은 $0$이다. 남는 것은 $j=k$와 $j=k-1$의 두 항이다. $j=k$의 계수가 $\alpha_k$이고, $j=k-1$에서는 $xp_{k-1}=p_k+(\text{차수 }k-1\text{ 이하})$이므로 $\langle p_k,xp_{k-1}\rangle=\langle p_k,p_k\rangle$이어서 계수가 $\beta_k$이다.
\end{proof}

이 증명은 27장의 Lanczos 과정에서 대칭 행렬의 Arnoldi 관계가 세 항으로 줄어든 것과 같은 구조이다. 실제로 곱셈 연산자 $f\mapsto xf$는 이 내적에 대해 대칭이고, 직교다항식은 그 연산자의 Krylov 부분공간 $\operatorname{span}\{1,x,\dots,x^k\}$의 직교 기저이다.

\begin{proposition}
\label{prop:orthopoly-roots}
$p_n$의 근은 모두 실수이고 서로 다르며 $(a,b)$ 안에 있다.
\end{proposition}

\begin{proof}
$p_n$이 $(a,b)$ 안에서 부호를 바꾸는 점, 곧 홀수 중복도의 근을 $y_1,\dots,y_m$이라 하자. $m<n$이라 가정하면 $q(x)=\prod_j(x-y_j)$는 차수가 $n$보다 낮으므로 $\langle p_n,q\rangle=0$이다. 그런데 $p_nq$는 $(a,b)$에서 부호를 바꾸지 않고 $0$이 아니므로 $\int p_nq\,\omega\ne0$이다. 모순이므로 $m=n$이고, 차수가 $n$이므로 이것이 근의 전부이며 모두 단순근이다.
\end{proof}

\section{Gauss 구적}
\label{sec:gauss}

\subsection{정확도와 양의 무게}

노드를 마음대로 고를 수 있다면 $n$개의 노드와 $n$개의 무게, 모두 $2n$개의 자유도로 차수 $2n-1$까지 맞출 수 있으리라 기대할 수 있다. 그 기대는 노드를 직교다항식의 근으로 둘 때 정확히 이루어진다\cite{gautschi2004orthogonal}.

\begin{theorem}[Gauss]
\label{thm:gauss}
$p_n$의 근 $x_1,\dots,x_n$을 노드로 하는 보간형 구적은 차수 $2n-1$ 이하의 모든 다항식 $f$에 대해 $\int f\omega=\sum_iw_if(x_i)$를 만족한다. 또 무게는 모두 양수이고, 차수 $2n$까지 정확한 $n$점 규칙은 없다.
\end{theorem}

\begin{proof}
차수 $2n-1$ 이하의 $f$를 $p_n$으로 나누어 $f=qp_n+r$로 쓰면 $q$와 $r$은 차수가 $n-1$ 이하이다. $\int qp_n\omega=0$이고 노드에서 $p_n=0$이므로 $f(x_i)=r(x_i)$이다. 보간형 구적은 차수 $n-1$ 이하의 $r$을 정확히 적분하므로 $\int f\omega=\int r\omega=\sum_iw_ir(x_i)=\sum_iw_if(x_i)$이다.

무게가 양수인 것은 $L_j^2$을 적분해 보면 안다. $L_j^2$은 차수 $2n-2$이므로 규칙이 정확하고, $0<\int L_j^2\omega=\sum_iw_iL_j(x_i)^2=w_j$이다. 마지막으로 $f(x)=\prod_i(x-x_i)^2$은 차수 $2n$이고 노드에서 $0$이지만 적분은 양수이므로 어떤 $n$점 규칙도 이것을 정확히 적분하지 못한다.
\end{proof}

무게가 양수이므로 \cref{sec:interpolatory-quadrature}의 조건수가 $1$이고, Weierstrass 정리와 합치면 Gauss 구적은 $n\to\infty$에서 모든 연속함수에 대해 수렴한다. Newton--Cotes와 정반대이다. 오차는 $f\in C^{2n}$에서 $f^{(2n)}(\xi)/(2n)!\cdot\int p_n^2\omega$이다.

\begin{aside}[Gauss의 구적]
Gauss는 1814년 괴팅겐 학회에서 이 구적법을 발표했는데, 그때는 직교다항식이라는 개념이 없었다. Gauss는 연분수와 초기하급수를 써서 노드가 Legendre 다항식의 근이 된다는 것을 이끌어 냈고 손으로 일곱 점 규칙까지의 노드와 무게를 열여섯 자리로 계산했다. 직교다항식으로 정리된 오늘날의 이해는 Jacobi와 Christoffel, Stieltjes를 거쳐 19세기 후반에 만들어졌다.
\end{aside}

\subsection{Golub--Welsch 방법}

\cref{thm:three-term}의 점화식을 정규화한 직교다항식 $\hat p_k=p_k/\|p_k\|$로 쓰면 $x\hat p_k=\sqrt{\beta_{k+1}}\hat p_{k+1}+\alpha_k\hat p_k+\sqrt{\beta_k}\hat p_{k-1}$이다. 벡터 $\hat p(x)=(\hat p_0(x),\dots,\hat p_{n-1}(x))^{\mathsf T}$와 대각이 $\alpha_k$, 부대각이 $\sqrt{\beta_k}$인 대칭 삼중대각 행렬 $J_n$을 쓰면 이것은
\[
J_n\hat p(x)=x\hat p(x)-\sqrt{\beta_n}\,\hat p_n(x)e_n
\]
이다. $x=x_i$가 $p_n$의 근이면 마지막 항이 사라지므로 $x_i$는 $J_n$의 고유값이고 $\hat p(x_i)$는 그 고유벡터이다. 이 $J_n$을 \term{Jacobi 행렬}{Jacobi matrix}이라 한다.

\begin{proposition}[Golub--Welsch]
\label{prop:golub-welsch}
$J_n$의 고유값 $x_i$와 단위 고유벡터 $v_i$에 대해 Gauss 무게는 $w_i=\mu_0(v_i)_1^2$이다. 여기서 $\mu_0=\int\omega$이다.
\end{proposition}

\begin{proof}
$P_{ki}=\hat p_k(x_i)$인 $n\times n$ 행렬 $P$를 두자. Gauss 구적은 차수 $2n-2$ 이하인 $\hat p_k\hat p_l$을 정확히 적분하므로 $\sum_iw_i\hat p_k(x_i)\hat p_l(x_i)=\delta_{kl}$, 곧 $PWP^{\mathsf T}=I$이다($W=\operatorname{diag}(w_i)$). $P$는 정칙이므로 $W=(P^{\mathsf T}P)^{-1}$이다. 그런데 $P$의 열은 대칭 행렬 $J_n$의 서로 다른 고유값에 속하는 고유벡터이므로 서로 직교하고, $P^{\mathsf T}P$는 대각 행렬 $\operatorname{diag}(\|\hat p(x_i)\|^2)$이다. 따라서 $w_i=1/\|\hat p(x_i)\|^2$이다. 단위 고유벡터는 $v_i=\hat p(x_i)/\|\hat p(x_i)\|$이고 첫 성분은 $\hat p_0=1/\sqrt{\mu_0}$을 $\|\hat p(x_i)\|$로 나눈 것이므로 $(v_i)_1^2=1/(\mu_0\|\hat p(x_i)\|^2)=w_i/\mu_0$이다.
\end{proof}

이 정리에 따라 Gauss 규칙은 23장의 대칭 삼중대각 고유값 문제 하나로 계산된다\cite{golub1969welsch}. 그러나 무게를 고유벡터의 첫 성분의 제곱에서 그대로 얻으면 정확도를 잃는다. 양 끝 근처의 노드에 대한 무게는 아주 작은데 고유벡터의 작은 성분은 절대오차로만 정확하므로 그 제곱의 상대오차가 커진다. 그래서 실제 라이브러리는 고유값으로 얻은 노드를 Newton 걸음으로 다듬고 무게는 $p_n'$을 쓰는 공식으로 다시 계산하는데, 그래도 무게의 마지막 자리들을 지키기는 쉽지 않다. $n=100$의 Gauss--Legendre 규칙의 참값을 Python의 \code{decimal}과 $60$자리 Newton 반복으로 구해 견주면, numpy의 \code{leggauss}는 노드가 $0.7$ulp 안이지만 무게는 최대 $6.3\times10^4$ulp 틀리고 scipy의 \code{roots\_legendre}는 노드가 $8.2$ulp, 무게가 $8.4\times10^4$ulp 틀린다. $n=5$에서도 무게가 numpy는 $11$ulp, scipy는 $17$ulp 틀린다. Legendre 다항식처럼 점화식의 계수를 알면 Newton 반복이 더 정확하다. 근의 점근 근사에서 출발해 $p_n$과 $p_n'$을 점화식으로 계산하며 반복한다. $n$이 아주 크면 반복 없이 점근 전개만으로 노드와 무게를 $O(n)$에 구하는 방법도 있다\cite{hale2013legendre,bogaert2014}.

\section{Clenshaw--Curtis 규칙}
\label{sec:clenshaw-curtis}

29장의 Chebyshev 점 $x_j=\cos(j\pi/n)$를 노드로 하는 보간형 구적이 \term{Clenshaw--Curtis 규칙}{Clenshaw--Curtis rule}이다\cite{clenshaw1960}. 정확도의 차수는 적어도 $n$($n$이 짝수면 $n+1$)으로 같은 개수의 점을 쓰는 Gauss 구적의 약 절반이지만, 무게가 모두 양수이고 FFT로 $O(n\log n)$에 계산된다\cite{waldvogel2006}. 실제 정확도가 놀랍다. Trefethen은 대부분의 함수에서 Clenshaw--Curtis가 Gauss 구적에 거의 뒤지지 않음을 보였다\cite{trefethen2008gauss}. 정확도의 차수는 다항식을 정확히 적분하는지만 정할 뿐이고 일반적인 함수의 오차는 Chebyshev 계수가 줄어드는 속도로 정해지기 때문이다.

$n+1$개의 점으로 $[-1,1]$의 세 함수를 적분하면 \cref{tab:gauss-cc}와 같다. $e^x$처럼 전 평면에서 해석적인 함수에서는 Gauss가 압도적이지만, 극점이 실수축 가까이 있는 Runge 함수에서는 두 배에서 다섯 배쯤, 매끄러움이 유한한 $|x|^3$에서는 거의 같다.

\begin{table}[t]
\begin{sdtable}{colspec={ccccccc}, row{1-2}={font=\small\bfseries}}
\toprule
$n$ & \SetCell[c=2]{c} $e^x$ & & \SetCell[c=2]{c} $1/(1+25x^2)$ & & \SetCell[c=2]{c} $|x|^3$ & \\
 & Gauss & CC & Gauss & CC & Gauss & CC \\
\midrule
$4$ & $8.2\times10^{-10}$ & $2.7\times10^{-5}$ & $1.6\times10^{-1}$ & $3.3\times10^{-1}$ & $2.1\times10^{-3}$ & $1.0\times10^{-2}$ \\
$8$ & $8.9\times10^{-16}$ & $2.0\times10^{-11}$ & $2.9\times10^{-2}$ & $5.5\times10^{-2}$ & $2.1\times10^{-4}$ & $4.4\times10^{-4}$ \\
$16$ & $0$ & $4.4\times10^{-16}$ & $1.2\times10^{-3}$ & $2.2\times10^{-3}$ & $1.8\times10^{-5}$ & $2.5\times10^{-5}$ \\
$32$ & $8.9\times10^{-16}$ & $0$ & $2.1\times10^{-6}$ & $3.9\times10^{-6}$ & $1.3\times10^{-6}$ & $1.6\times10^{-6}$ \\
$64$ & $2.7\times10^{-15}$ & $4.4\times10^{-16}$ & $6.2\times10^{-12}$ & $2.9\times10^{-11}$ & $8.8\times10^{-8}$ & $9.7\times10^{-8}$ \\
\bottomrule
\end{sdtable}
\caption{점 $n+1$개의 Gauss--Legendre와 Clenshaw--Curtis(CC) 규칙의 절대오차. numpy로 계산했다.}
\label{tab:gauss-cc}
\end{table}

\section{오차의 어림: Gauss--Kronrod}
\label{sec:kronrod}

적분 루틴은 값과 함께 오차의 어림을 내야 한다. 가장 단순하게는 서로 다른 두 규칙의 차이를 본다. 다만 두 규칙이 노드를 공유하지 않으면 함수 호출이 두 배로 든다. Kronrod는 $n$점 Gauss 규칙의 노드에 $n+1$개의 노드를 더해 원래 노드를 그대로 쓰면서 정확도의 차수를 적어도 $3n+1$($n$이 홀수면 $3n+2$)로 올리는 규칙을 만들었다\cite{kronrod1965}. 더한 노드는 Stieltjes 다항식이라 부르는 특별한 다항식의 근이고, Laurie는 그 노드와 무게를 Golub--Welsch와 비슷한 고유값 문제로 계산하는 방법을 내놓았다\cite{laurie1997kronrod}. 이렇게 만든 \term{Gauss--Kronrod 쌍}{Gauss--Kronrod pair}은 $2n+1$번의 호출로 두 값 $G_n$과 $K_{2n+1}$을 주고, 더 정확한 $K$를 결과로 쓰며 $|K-G|$를 오차의 어림으로 쓴다.

$|K-G|$는 사실 덜 정확한 $G$의 오차의 어림이므로 $K$의 오차에 대해서는 지나치게 비관적이다. Piessens 등의 QUADPACK은 경험적인 보정 $\mathrm{err}=I_{|f|}\min\bigl(1,(200|K-G|/I_{|f|})^{3/2}\bigr)$에 반올림의 바닥 $50\varepsilon\int|f|$를 더한 $\max(50\varepsilon\int|f|,\ \mathrm{err})$을 쓴다\cite{piessens1983quadpack}. 여기서 $I_{|f|}$는 $|f-\bar f|$를 같은 규칙으로 적분한 값이다. 오차가 작을수록 $3/2$제곱이 어림을 더 줄여 실제 오차에 가깝게 한다. 그래도 어림은 대개 실제 오차보다 크다. scipy의 \code{quad}로 $[0,1]$의 $e^x$를 적분하면 21점 규칙 하나로 끝나고 실제 오차가 $2.2\times10^{-16}$, 어림이 $1.9\times10^{-14}$이다. 이 어림은 바로 반올림의 바닥 $50\varepsilon(e-1)$이다. $[0,1]$의 $\cos100x$는 구간 $8$개로 나누어 실제 오차 $6.6\times10^{-16}$, 어림 $5.9\times10^{-10}$이다.

어림이 실제 오차보다 작으면 더 위험하다. 오차의 어림은 함수를 본 점에서만 계산되므로 처음 21점이 모두 함수의 중요한 부분을 놓치면 규칙은 틀린 값을 정확하다고 판정한다. $f(x)=4^{-\alpha}/\bigl((x-\pi/4)^2+16^{-\alpha}\bigr)$은 $x=\pi/4$에 폭이 $4^{-\alpha}$인 봉우리가 있고 $[0,1]$의 적분이 거의 $\pi$이다. scipy의 \code{quad}는 $\alpha=19$에서 $3.5\times10^{-9}$를, $\alpha=20$에서 $8.6\times10^{-10}$을 내면서 경고 없이 수렴했다고 알리고, 오차의 어림은 각각 $5.2\times10^{-9}$, $1.3\times10^{-9}$이다. 함수를 상자로만 보는 적분 루틴은 이 한계를 피할 수 없다. 봉우리의 위치를 알면 거기서 구간을 나누어 주어야 한다.

\section{stardust의 구적 규칙}
\label{sec:quadrature-stardust}

\begin{instardust}[\code{docs/design/integrate.md} §3.7]
stardust의 \code{fixed\_quad}는 $n$점 Gauss--Legendre 규칙이다. 노드와 무게는 올바르게 반올림한 binary64이다. $n\le100$은 Arb로 Legendre 다항식의 근을 $2^{-110}$까지 정해 반올림한 표를 쓰고, 반올림의 경계에 걸리는 값이 있으면 표를 만드는 도구가 멈춘다. $n>100$은 점마다 점근 근사에서 시작해 binary64와 double-double의 Newton 반복으로 노드를, 그 노드에서 무게를 계산하고, 반올림의 경계에 가까우면 정확한 정수 산술로 정한다. 설계 문서는 scipy의 \code{roots\_legendre}가 $n=100$에서 노드를 $8$ulp, 무게를 $83682$ulp까지 틀린다고 적었는데 위에서 \code{decimal}로 독립적으로 잰 값과 같다. \code{newton\_cotes}의 무게와 오차 계수는 정확한 유리수로 계산해 한 번 반올림한다.
\end{instardust}

\begin{instardust}[\code{docs/design/integrate.md} §3.1]
\code{quad}는 QUADPACK의 일곱 루틴을 scipy의 C 판에서 문장 단위로 옮겼다. Kronrod 규칙의 상수는 QUADPACK의 $10$진 리터럴 그대로이다. 34장의 근 찾기처럼 scipy의 C 코드는 clang이 곱과 합을 FMA로 합쳐 빌드된다. 문장마다 옮겼을 때 합쳐지는 곳은 Kronrod 규칙의 누적 넷뿐이어서 그 자리에 \code{fma}를 적었다. 그 결과 시험 문제 $106$개 모두에서 함수를 부른 점, 값, 오차의 어림, 호출 수가 scipy와 비트까지 같다. 위의 좁은 봉우리처럼 어림이 틀리는 경우도 scipy와 같게 두고, 문서에 봉우리에서 구간을 나누는 법을 적었다.
\end{instardust}

\section{맺음}

구적 규칙의 질은 노드를 어디에 두느냐와 무게가 양수냐로 정해진다. 고른 간격의 노드는 차수를 올리면 음수 무게 때문에 불안정해지고, 직교다항식의 근에 둔 노드는 차수를 두 배로 올리면서도 무게를 양수로 지킨다. 그러나 차수가 전부는 아니다. 주기 함수에서는 사다리꼴 규칙이, 대부분의 함수에서는 Clenshaw--Curtis 규칙이 Gauss 구적과 맞먹는다. 다음 장에서는 규칙 하나를 고르는 데서 나아가 구간을 적응적으로 나누고 수열을 가속하고 변수를 바꾸어 특이점과 무한 구간을 다루는 방법을 본다.

수치 적분의 고전적인 교과서는 Davis와 Rabinowitz의 책\cite{davis1984methods}이고, 직교다항식과 Gauss 구적의 계산은 Gautschi의 책\cite{gautschi2004orthogonal}이 자세하다.

\chapterreferences
